\documentclass[10pt,a4paper]{amsart}
\usepackage[margin=19mm]{geometry}
\usepackage{amsmath,amssymb,mathtools}
\usepackage[scr=rsfs,cal=euler]{mathalfa}
\usepackage{xfrac}
\usepackage[unicode,hidelinks,bookmarks=false]{hyperref}
\newcommand{\ie}{i.e.\ }
\newcommand{\cf}{cf.\ }
\newcommand{\resp}{respectively\ }
\newcommand{\Subsection}{\subsection}
\providecommand{\nobreakdash}{\nobreak\hbox{-}}
\setlength{\emergencystretch}{2em}
\multlinegap0pt
\usepackage{amssymb}
\usepackage{xcolor}
\newtheorem{corollary}{Corollary}
\newtheorem{proposition}{Proposition}
\newtheorem{theorem}{Theorem}
\title{A Generalization of $\Delta$U Rings}
\author{Peter Danchev, Omid Hasanzadeh, Ahmad Moussavi, Mehrdad Esfandiar}
\date{Source: arXiv:2605.22700v1 (CC BY 4.0)}
\begin{document}
\maketitle

\noindent\emph{Source and licence.} A Generalization of $\Delta$U Rings. Version arXiv:2605.22700v1 (CC BY 4.0), distributed by arXiv under the Creative Commons Attribution 4.0 International licence, http://creativecommons.org/licenses/by/4.0/, as recorded in the arXiv licence metadata for this exact version and shown on the abstract page https://arxiv.org/abs/2605.22700v1. Full licence notice: \texttt{licenses/arxiv-2605.22700-CC-BY-4.0.txt}.\par\smallskip

\noindent\textit{Editorial scope.} Counted unit: the article's theorem fin2 (source0.tex lines 799--801). The article's complete original proof of it is delivered with it (source0.tex lines 803--828, one \texttt{proof} environment of 26 lines). No mathematics is written, repaired or re-derived: the statement and the proof are the article's own lines, in the article's own notation, and no retained line is substituted or re-flowed.  Retained uncounted as the source-local context the proof needs: lines 365--367; lines 552--562; lines 698--712.  Nothing was omitted from any retained span, so the package carries no omission record.  No retained line contains a citation, so the wrapper carries no bibliography and no bibliography entry.  No retained line contains a figure environment, an image-inclusion command or drawing code, so no image is delivered and the package declares no asset dependency.  Editorial formatting only, in addition: an AMS \texttt{amsart} preamble that restates the article's own declarations and macros used by the retained lines, namely the environments \texttt{corollary}, \texttt{proposition}, \texttt{theorem} on separate counters, together with the standard packages those lines need.  The retained environments print in this document as Corollary 1, Proposition 1, Theorem 1 in order of appearance, whereas the article prints the same items with the numbering of its own full document.  The complete native source of the article as distributed in the arXiv e-print for this exact version is bundled unchanged as \texttt{sources/201152/source0.tex}.
\begin{corollary}\label{2}
A ring \(R\) is W$\Delta$U if, and only if, \(R/J(R)\) is W$\Delta$U.
\end{corollary}

\begin{proposition}\label{division}
Let \( R \) be a ring. Then, the following four conditions are valid:
	
(1) A division ring \( R \) is W$\Delta$U if, and only if, \( R \cong \mathbb{Z}_2 \) or \( \mathbb{Z}_3 \).
	
(2) A local ring \( R \) is W$\Delta$U if, and only if, either \( R/J(R) \cong \mathbb{Z}_2 \) or \(R/J(R) \cong  \mathbb{Z}_3 \).
	
(3) A semi-simple ring $R$ is W$\Delta$U if, and only if, $R \cong \bigoplus_{i=1}^n R_i$, where, for just only one index $i$, $R_i\cong \mathbb{Z}_3 \text{ and } R_j\cong \mathbb{Z}_2$ for any other index $j\neq i$.
	
(4) A semi-local ring $R$ is W$\Delta$U if, and only if, $R/J(R) \cong \bigoplus_{i=1}^m R_i$, where, for just only one index $i$, $R_i\cong \mathbb{Z}_3 \text{ and } R_j\cong \mathbb{Z}_2$ for any other index $j\neq i$.	
\end{proposition}

\begin{theorem}\label{main}
Suppose \( R \) is a ring. Then, the following four conditions are equivalent:
	
(1) \( R \) is exchange W$\Delta$U.
	
(2) \( R \) is clean W$\Delta$U.
	
(3) \( R \) is weakly $\Delta$-clean.
	
(4) All idempotents of \( R \) lift modulo \( J(R) \), and either
\[
	R / J(R) \cong B, \quad \text{or} \quad R / J(R) \cong \mathbb{Z}_3, \quad \text{or} \quad R / J(R) \cong B \times \mathbb{Z}_3,
\]
where \( B \) is a Boolean ring.
\end{theorem}

\begin{theorem}\label{fin2}
A ring \( R \) is weakly clean W\(\Delta U\) with \( 2 \in U(R) \) if, and only if, \( R / J(R) \cong \mathbb{Z}_3 \).
\end{theorem}

\begin{proof}
Since the sufficiency has already been established in Theorem \ref{main}, we now turn our attention to the necessity. We manage to demonstrate that \( R \) is indecomposable; that is, the only idempotents in \( R \) are \( \{0, 1\} \).

To this intention, consider an arbitrary idempotent \( e \in R \). Note that the element \( 2e - 1 \) is a unit in \( R \), and hence it can be written as either \( 2e - 1 = 1 + d \) or \( 2e - 1 = -1 + d \), where \( d \in \Delta(R) \). Consequently, we obtain \( 2e = 2 + d \) or \( 2e = d \), yielding that \( e = 1 + \frac{d}{2} \) or \( e = \frac{d}{2} \). In both cases, \( e \in U(R) \) or \( e \in \Delta(R) \), asserting that \( e = 1 \) or \( e = 0 \).

That is why, there are precisely six possibilities for any element \( r \in R \) expressed in terms of \( d \in \Delta(R) \) as follows:

\medskip

\begin{itemize}
	\item \( r = 1 + d + 1 = 2 + d \in U(R) + \Delta(R) \subseteq U(R) \),
	\item \( r = 1 + d - 1 = d \in \Delta(R) \),
	\item \( r = 1 + d \in 1 + \Delta(R) \subseteq U(R) \),
	\item \( r = -1 + d - 1 = -2 + d \in U(R) \),
	\item \( r = -1 + d + 1 = d \in \Delta(R) \),
	\item \( r = -1 + d \in U(R) \).
\end{itemize}

\medskip

Hence, each element of \( R \) is either a unit or lies in \( \Delta(R) \). This enables us that \( R \) is a local ring, and so \( R/J(R) \) is a division ring.

To verify this, let \( x \in R \) be such that \( x \not\in J(R) \). Then, there exists \( y \in R \) such that \( 1 - xy \not\in U(R) \) enabling us that \( 1 - xy \in \Delta(R) \), and thus \( xy \in 1 + \Delta(R) \subseteq U(R) \). Also, for some \( h \in R \), the element \( hx \) is also a unit, guaranteeing that \( x \) itself must be a unit too.

Moreover, since \( R/J(R) \) is a W\(\Delta\)U-ring in agreement with Corollary \ref{2}, as well as it is simultaneously a division ring, it follows from Proposition \ref{division} that \( R/J(R) \cong \mathbb{Z}_3 \), as needed.
\end{proof}

\end{document}
